\documentclass[11pt]{article}
\usepackage[margin=1in]{geometry}
\usepackage{lmodern}
\usepackage{hyperref}
\usepackage{enumitem}
\usepackage{amsmath,amssymb,amsthm,mathtools}
\usepackage[T1]{fontenc}
\usepackage[utf8]{inputenc}
\title{\textbf{Theory-mode report: Anchored random clusters and SLE excursions}}
\author{Generated from Scholar Inbox digest}
\date{Digest date: 2026-05-07}
\begin{document}
\maketitle
\section*{Paper information}
\begin{itemize}[leftmargin=1.5em]
\item \textbf{Title:} Anchored random clusters and SLE excursions
\item \textbf{Authors:} Valentino Foit, Jesper Lykke Jacobsen, Hjalmar Rosengren
\item \textbf{arXiv:} \href{https://arxiv.org/abs/2605.04395}{2605.04395}
\end{itemize}
\section*{Executive summary}
This paper is a pedagogical but nontrivial review of conformal-field-theory methods for computing Schramm--Loewner evolution observables in the upper half-plane. The central theme is that a range of probabilistic quantities associated with SLE curves and critical random clusters can be represented as bulk-boundary correlation functions involving degenerate boundary fields. Once that representation is in place, BPZ null-vector equations turn the probability question into a differential-equation problem.

The paper revisits standard examples such as Schramm's left-passage probability and SLE Green's functions, and then explains how the same machinery yields generalized densities of \emph{anchored} critical percolation clusters, together with new formulas for densities of pivotal points between critical Fortuin--Kasteleyn clusters. So the paper is partly expository, but it also pushes the anchored-cluster story beyond the classical percolation case.
\section*{General setup}
The setting is chordal SLE in the upper half-plane $\mathbb H$, coupled conceptually to two-dimensional critical lattice models through conformal invariance and CFT. Certain observables---for example, whether the SLE trace passes to the left of a marked point, or the local density of a critical cluster anchored at boundary points---are encoded as correlation functions of bulk operators with special boundary condition changing fields.

The decisive input is that some boundary operators are degenerate at level two. Their insertion forces the correlator to satisfy second-order BPZ differential equations. The probabilistic observable is then recovered from the solution selected by conformal covariance and the relevant boundary conditions.
\section*{Main examples treated in the paper}
\subsection*{1. Schramm's left-passage probability}
The simplest example is the probability that a chordal SLE trace passes to the left of a given bulk point. In the CFT language, this becomes a four-point bulk-boundary correlator with one degenerate boundary insertion. The BPZ equation reduces the observable to an ordinary differential equation in the cross-ratio-like angular variable. Solving that equation with the correct boundary behavior reproduces Schramm's formula.

The point here is not just to re-derive a known result. It sets up the full dictionary between SLE observables and degenerate conformal blocks.
\subsection*{2. SLE Green's functions}
The Green's function measures the normalized probability that the SLE curve comes close to a marked bulk point. Again the observable satisfies a BPZ equation, and the CFT method recovers the known power-law dependence and angular structure. This is the prototype for local density observables attached to random curves.
\subsection*{3. Anchored critical clusters}
The more distinctive topic is the density of clusters constrained to touch prescribed boundary anchor points. The paper explains how formulas first obtained by Kleban, Simmons, and Ziff arise naturally from CFT correlators. In this language, the anchored cluster density is not an ad hoc lattice formula but a conformal observable with explicit operator content.

One of the useful outputs is a unified derivation of one-anchor, two-anchor, and related densities using the same degenerate-field framework.
\subsection*{4. New formulas for pivotal-point densities in FK clusters}
The genuinely new piece advertised in the abstract is the extension to densities of pivotal points between critical Fortuin--Kasteleyn clusters. These are more delicate observables because they probe the geometry of connectivity changes. The same CFT/SLE machinery still applies, but the relevant correlators and fusion channels are subtler.
\section*{Why degenerate operators matter}
In boundary CFT, inserting a degenerate primary creates a null descendant, and decoupling of that null state yields a differential equation for the correlator. For SLE this is exactly what one wants: a stochastic geometric probability is converted into a deterministic PDE/ODE problem constrained by conformal invariance.

So the logic is:
\begin{itemize}[leftmargin=1.5em]
\item identify the probabilistic event or density,
\item represent it as a bulk-boundary correlator,
\item use a degenerate boundary operator to derive BPZ equations,
\item solve the equations subject to the correct asymptotics and boundary conditions.
\end{itemize}
This is the common backbone behind left-passage, Green's functions, and anchored cluster densities.
\section*{Structure of the paper}
From the source structure, the paper is organized as follows.
\begin{itemize}[leftmargin=1.5em]
\item \textbf{Introduction:} motivation from SLE observables and anchored clusters.
\item \textbf{Loop models, SLE, and CFT:} the probabilistic and conformal background.
\item \textbf{A section on CFT technology:} likely introducing the relevant boundary operators, null-state equations, and conformal blocks.
\item \textbf{A section on concrete probabilities/densities:} where left-passage, Green's functions, and anchored-cluster formulas are derived.
\item \textbf{Final section:} extension to pivotal FK-cluster densities and concluding interpretation.
\end{itemize}
So despite the pedagogical tone, it is really a worked guide to how one computes these observables in practice.
\section*{Conceptual significance}
The paper is valuable because it makes the SLE/CFT dictionary operational. Many people know at a slogan level that ``SLE observables correspond to CFT correlators.'' This paper actually walks through the mechanism for a family of concrete observables and shows how the same formalism adapts to anchored cluster densities.

For statistical mechanics, the anchored-cluster viewpoint is especially nice because it connects continuum conformal methods to visually intuitive lattice events. For probability, it clarifies why these densities satisfy rigid differential constraints and why hypergeometric-type formulas keep appearing.
\section*{Limitations}
This is not a general theorem proving new universality from first principles. The method assumes the conformal/SLE/CFT framework and then computes observables within it. Also, because the paper is partly pedagogical, some of its value lies in synthesis and derivation rather than in a single sharp theorem. Still, the new FK pivotal-point formulas mean it is not merely a survey.
\section*{Bottom line}
A useful and readable paper if you care about the interface between SLE, conformal field theory, and critical cluster geometry. The main lesson is that anchored random-cluster densities are naturally encoded as bulk-boundary correlators with degenerate insertions, and once you accept that, the observable formulas follow from BPZ equations plus conformal boundary conditions. The new pivotal-point formulas make the review worth attention even for people who already know the classical left-passage story.
\end{document}
